\documentclass[12pt]{article}
\usepackage[utf8]{inputenc}
\usepackage[T1]{fontenc}
\usepackage{hyperref}
\title{Sanga Language Resources for K-12 Curriculum Development}
\author{Compiled by Claude Sonnet 5.5 (claude-sonnet-5-5) for Andr\'as Kornai}
\date{}
\begin{document}
\maketitle

\section{Introduction}
Sanga (ISO 639-3 xsn; Glottocode sang1329) is an East Kainji language of Nigeria, within the Atlantic-Congo/Niger-Congo family. The sources found place it in about 20 villages in Toro LGA, Bauchi State, and Jema'a LGA, Kaduna State. Wikipedia gives roughly 20,000 speakers (a 2000 figure). Joshua Project gives about 41,000 Sanga people. Alternate names include Asanga, Isanga, Sangawa and Sangauchi. Only a few resources turned up in the searches. No school materials, primers or textbooks were found, so this report is short.

\section{Primary Educational Resources}
No primers, textbooks or school programmes for Sanga were found in the search results.

\section{Linguistic and Reference Materials}
\begin{itemize}
\item Wikipedia, ``Sanga language (Nigeria)'': a short overview covering classification, location and speaker numbers. \url{https://en.wikipedia.org/wiki/Sanga_language_(Nigeria)}
\item Lexvo/Finto entry for xsn, a language-code record only. \url{https://finto.fi/lexvo/en/page/xsn}
\end{itemize}

\section{Religious and Community Resources}
\begin{itemize}
\item Joshua Project, Sanga in Nigeria: people-group profile. It reports that a Bible translation is in progress and that the community is mostly Christian with some Muslims. \url{https://joshuaproject.net/people_groups/14721/NI}
\item Joshua Project, Sanga language page. \url{https://www.joshuaproject.net/index.php/languages/xsn}
\end{itemize}

\section{Digital and Audio Resources}
\begin{itemize}
\item Global Recordings Network, Sanga: audio recordings for evangelism and basic Bible teaching. \url{https://globalrecordings.net/en/language/xsn}
\end{itemize}
Joshua Project also reports that the Jesus Film is available in Sanga.

\section{Curriculum Development Context}
No evidence was found of Sanga being taught in schools. Any curriculum would have to start from new documentation: orthography development, word lists and community-produced materials. The audio and Bible-translation work noted above could be a starting point. Contact with the Sanga community and the translation project would be needed. The speaker counts found in the searches are rough and the sources disagree.
\end{document}